% =================================================
% TikZ-рисунки к листу «Ломаная и многоугольник»
% =================================================

\tikzset{
  geom edge/.style={draw=Secondary,line width=1.45pt,line cap=round,line join=round},
  geom side/.style={draw=Primary,line width=1.55pt,line cap=round,line join=round},
  geom diagonal/.style={draw=Accent,line width=1.15pt},
  geom auxiliary/.style={draw=GrayLine,line width=0.75pt,densely dashed},
  geom point/.style={circle,fill=Accent,draw=white,line width=0.35pt,inner sep=2.15pt},
  geom label/.style={font=\small\bfseries,text=GrayText},
  geom caption/.style={font=\small\sffamily,text=GrayText,align=center},
  geom note/.style={font=\small\bfseries\sffamily,text=Secondary,align=center},
  geom card/.style={draw=GrayLine,fill=white,rounded corners=1.5mm,line width=0.65pt},
  geom region one/.style={fill=PrimaryLight,draw=none},
  geom region two/.style={fill=SecondaryLight,draw=none},
  geom region three/.style={fill=AccentLight,draw=none}
}

\newcommand{\gpoint}[2]{%
  \node[geom point] at (#1) {};
  \node[geom label,#2] at (#1) {};%
}

\newcommand{\openbrokenfigure}{%
  \begin{center}
    \begin{tikzpicture}[x=1cm,y=1cm]
      \coordinate (A) at (-5.0,-0.5);
      \coordinate (B) at (-2.4,1.2);
      \coordinate (C) at (0.0,-0.25);
      \coordinate (D) at (2.4,1.15);
      \coordinate (E) at (5.0,-0.55);
      \draw[geom edge] (A)--(B)--(C)--(D)--(E);
      \foreach \P/\pos in {A/below left,B/above,C/below,D/above,E/below right}{
        \node[geom point] at (\P) {};
        \node[geom label,\pos=3pt] at (\P) {$\P$};
      }
      \node[geom caption] at (0,-1.25)
        {Ломаная $ABCDE$: вершины $A,B,C,D,E$; звенья $AB,BC,CD,DE$};
    \end{tikzpicture}
  \end{center}
}

\newcommand{\openclosedfigure}{%
  \begin{center}
    \begin{tikzpicture}[x=0.9cm,y=0.9cm]
      \node[geom card,minimum width=7.0cm,minimum height=4.4cm] at (-4.3,0) {};
      \coordinate (A) at (-6.7,-0.6); \coordinate (B) at (-5.5,1.0);
      \coordinate (C) at (-3.9,-0.4); \coordinate (D) at (-2.2,0.9);
      \draw[geom edge] (A)--(B)--(C)--(D);
      \foreach \P/\pos in {A/below,B/above,C/below,D/above}{
        \node[geom point] at (\P) {}; \node[geom label,\pos=3pt] at (\P) {$\P$};}
      \node[geom note] at (-4.45,-1.65) {незамкнутая};

      \node[geom card,minimum width=7.0cm,minimum height=4.4cm] at (4.3,0) {};
      \coordinate (K) at (1.8,-0.65); \coordinate (L) at (2.6,1.05);
      \coordinate (M) at (4.55,1.3); \coordinate (N) at (6.7,-0.35);
      \coordinate (P) at (4.35,-1.05);
      \draw[geom edge] (K)--(L)--(M)--(N)--(P)--cycle;
      \foreach \Pnt/\lab/\pos in {K/K/below left,L/L/above left,M/M/above,N/N/right,P/P/below}{
        \node[geom point] at (\Pnt) {}; \node[geom label,\pos=3pt] at (\Pnt) {$\lab$};}
      \node[geom note] at (4.35,-1.65) {замкнутая};
    \end{tikzpicture}
  \end{center}
}

\newcommand{\polygonpartsfigure}{%
  \begin{center}
    \begin{tikzpicture}[x=1cm,y=1cm]
      \coordinate (A) at (-4.5,-1.25);
      \coordinate (B) at (-3.2,1.5);
      \coordinate (C) at (0.0,2.0);
      \coordinate (D) at (3.65,0.95);
      \coordinate (E) at (3.0,-1.6);
      \coordinate (F) at (-0.7,-2.0);
      \fill[PrimaryLight] (A)--(B)--(C)--(D)--(E)--(F)--cycle;
      \draw[geom side] (A)--(B)--(C)--(D)--(E)--(F)--cycle;
      \foreach \P/\pos in {A/left,B/above left,C/above,D/right,E/below right,F/below}{
        \node[geom point] at (\P) {}; \node[geom label,\pos=3pt] at (\P) {$\P$};}
      \node[geom note] at (0,-0.1) {шестиугольник $ABCDEF$};
      \node[geom caption,align=left] at (-5.5,1.9) {вершины};
      \draw[geom auxiliary] (-4.8,1.65)--(B);
      \node[geom caption,align=left] at (5.1,1.5) {стороны};
      \draw[geom auxiliary] (4.45,1.25)--($(C)!0.55!(D)$);
    \end{tikzpicture}
  \end{center}
}

\newcommand{\polygonnamesfigure}{%
  \begin{center}
    \begin{tikzpicture}[x=0.82cm,y=0.82cm]
      \foreach \x/\n/\label in {-6.0/3/треугольник,-2.0/4/четырёхугольник,2.0/5/пятиугольник,6.0/6/шестиугольник}{
        \ifnum\n=3
          \draw[geom side] (\x-1.0,-0.8)--(\x,1.0)--(\x+1.1,-0.8)--cycle;
        \fi
        \ifnum\n=4
          \draw[geom side] (\x-1.1,-0.75)--(\x-0.7,0.95)--(\x+1.0,0.75)--(\x+1.15,-0.8)--cycle;
        \fi
        \ifnum\n=5
          \draw[geom side] (\x-1.15,-0.65)--(\x-0.75,0.75)--(\x,1.15)--(\x+1.1,0.45)--(\x+0.85,-0.9)--cycle;
        \fi
        \ifnum\n=6
          \draw[geom side] (\x-1.15,-0.4)--(\x-0.8,0.75)--(\x,1.1)--(\x+1.0,0.65)--(\x+1.15,-0.45)--(\x,-1.0)--cycle;
        \fi
        \node[geom caption,text width=3.2cm] at (\x,-1.55) {\label};
      }
    \end{tikzpicture}
  \end{center}
}

\newcommand{\correctnamingfigure}{%
  \begin{center}
    \begin{tikzpicture}[x=1cm,y=1cm]
      \coordinate (A) at (-3.7,-1.15); \coordinate (B) at (-3.0,1.15);
      \coordinate (C) at (-0.3,1.8); \coordinate (D) at (3.2,0.7);
      \coordinate (E) at (2.35,-1.55); \coordinate (F) at (-0.8,-1.9);
      \draw[geom side] (A)--(B)--(C)--(D)--(E)--(F)--cycle;
      \foreach \P/\pos in {A/left,B/above left,C/above,D/right,E/below right,F/below}{
        \node[geom point] at (\P) {}; \node[geom label,\pos=3pt] at (\P) {$\P$};}
      \node[geom caption,text width=14cm] at (0,-2.55)
        {Называем вершины подряд по границе: $ABCDEF$, $BCDEFA$, $AFEDCB$ и т. п.};
    \end{tikzpicture}
  \end{center}
}

\newcommand{\diagonalfigure}{%
  \begin{center}
    \begin{tikzpicture}[x=1cm,y=1cm]
      \coordinate (A) at (-4.4,-1.35); \coordinate (B) at (-3.1,1.45);
      \coordinate (C) at (0.2,2.0); \coordinate (D) at (3.8,0.9);
      \coordinate (E) at (2.55,-1.65); \coordinate (F) at (-0.9,-2.0);
      \fill[PrimaryLight] (A)--(B)--(C)--(D)--(E)--(F)--cycle;
      \draw[geom side] (A)--(B)--(C)--(D)--(E)--(F)--cycle;
      \draw[geom diagonal] (A)--(C);
      \draw[geom diagonal] (A)--(D);
      \draw[geom diagonal] (A)--(E);
      \foreach \P/\pos in {A/left,B/above left,C/above,D/right,E/below right,F/below}{
        \node[geom point] at (\P) {}; \node[geom label,\pos=3pt] at (\P) {$\P$};}
      \node[geom note] at (0,-0.25) {диагонали из вершины $A$: $AC$, $AD$, $AE$};
    \end{tikzpicture}
  \end{center}
}

\newcommand{\pentagonsplitfigure}{%
  \begin{center}
    \begin{tikzpicture}[x=1cm,y=1cm]
      \coordinate (A) at (-4.3,-1.4); \coordinate (B) at (-3.2,1.2);
      \coordinate (C) at (0.0,2.0); \coordinate (D) at (3.8,0.65);
      \coordinate (E) at (2.25,-1.75);
      \fill[PrimaryLight] (A)--(B)--(C)--cycle;
      \fill[SecondaryLight] (A)--(C)--(D)--(E)--cycle;
      \draw[geom side] (A)--(B)--(C)--(D)--(E)--cycle;
      \draw[geom diagonal] (A)--(C);
      \foreach \P/\pos in {A/left,B/above left,C/above,D/right,E/below}{
        \node[geom point] at (\P) {}; \node[geom label,\pos=3pt] at (\P) {$\P$};}
      \node[geom note] at (-1.9,0.65) {$\triangle ABC$};
      \node[geom note] at (1.2,-0.35) {четырёхугольник $ACDE$};
    \end{tikzpicture}
  \end{center}
}

\newcommand{\quadrilateralregionsfigure}{%
  \begin{center}
    \begin{tikzpicture}[x=1cm,y=1cm]
      \coordinate (A) at (-4.4,-1.6); \coordinate (B) at (-3.3,1.55);
      \coordinate (C) at (3.6,1.2); \coordinate (D) at (4.0,-1.75);
      \coordinate (M) at ($(A)!0.5!(D)$);
      \path[name path=AC] (A)--(C);
      \path[name path=BM] (B)--(M);
      \path[name intersections={of=AC and BM,by=P}];
      \fill[PrimaryLight] (A)--(B)--(P)--cycle;
      \fill[SecondaryLight] (B)--(C)--(P)--cycle;
      \fill[AccentLight] (A)--(P)--(M)--cycle;
      \fill[GrayLight] (P)--(C)--(D)--(M)--cycle;
      \draw[geom side] (A)--(B)--(C)--(D)--cycle;
      \draw[geom diagonal] (A)--(C);
      \draw[geom diagonal] (B)--(M);
      \foreach \Pnt/\pos in {A/left,B/above left,C/above right,D/right,M/below,P/above}{
        \node[geom point] at (\Pnt) {}; \node[geom label,\pos=3pt] at (\Pnt) {$\Pnt$};}
    \end{tikzpicture}
  \end{center}
}

\newcommand{\studentrecognizefigure}{%
  \begin{center}
    \begin{tikzpicture}[x=0.82cm,y=0.82cm]
      \foreach \x/\lab in {-6/{1},-2/{2},2/{3},6/{4}}{
        \node[geom card,minimum width=3.5cm,minimum height=3.5cm] at (\x,0) {};
        \node[font=\bfseries\sffamily,text=Primary] at (\x,1.35) {\lab};
      }
      \draw[geom edge] (-7.2,-0.7)--(-6.35,0.55)--(-5.65,-0.35)--(-4.9,0.65);
      \draw[geom edge] (-3.15,-0.7)--(-2.8,0.65)--(-1.3,0.8)--(-0.85,-0.65)--cycle;
      \draw[geom edge] (0.8,-0.6)--(1.4,0.75)--(2.2,-0.45)--(3.1,0.65);
      \draw[geom edge] (4.85,-0.7)--(5.25,0.7)--(6.25,0.95)--(7.15,0.0)--(6.7,-0.85)--cycle;
      \draw[geom diagonal] (5.25,0.7)--(6.7,-0.85);
    \end{tikzpicture}
  \end{center}
}

\newcommand{\studentblankpentagon}{%
  \begin{center}
    \begin{tikzpicture}[x=1cm,y=1cm]
      \node[draw=GrayLine,line width=0.6pt,minimum width=14.8cm,minimum height=7.0cm] at (0,0) {};
      \node[geom caption] at (0,2.55) {Место для построения};
    \end{tikzpicture}
  \end{center}
}

\newcommand{\studentquadrilateralfigure}{%
  \begin{center}
    \begin{tikzpicture}[x=1cm,y=1cm]
      \coordinate (A) at (-4.5,-1.4); \coordinate (B) at (-3.4,1.4);
      \coordinate (C) at (3.5,1.15); \coordinate (D) at (4.1,-1.55);
      \coordinate (M) at ($(A)!0.5!(D)$);
      \draw[geom side] (A)--(B)--(C)--(D)--cycle;
      \foreach \P/\pos in {A/left,B/above left,C/above right,D/right,M/below}{
        \node[geom point] at (\P) {}; \node[geom label,\pos=3pt] at (\P) {$\P$};}
      \node[geom caption] at (0,-2.1) {Проведите $AC$ и $BM$. Точку их пересечения обозначьте $P$.};
    \end{tikzpicture}
  \end{center}
}

\newcommand{\hexagondiagonalsstudent}{%
  \begin{center}
    \begin{tikzpicture}[x=0.86cm,y=0.76cm]
      \coordinate (A) at (-3.9,0); \coordinate (B) at (-2.0,1.8);
      \coordinate (C) at (1.0,1.8); \coordinate (D) at (3.9,0);
      \coordinate (E) at (1.5,-1.8); \coordinate (F) at (-2.2,-1.8);
      \draw[geom side] (A)--(B)--(C)--(D)--(E)--(F)--cycle;
      \foreach \P/\pos in {A/left,B/above left,C/above,D/right,E/below,F/below}{
        \node[geom point] at (\P) {}; \node[geom label,\pos=3pt] at (\P) {$\P$};}
    \end{tikzpicture}
  \end{center}
}

\newcommand{\hexagondiagonalsanswer}{%
  \begin{center}
    \begin{tikzpicture}[x=1cm,y=1cm]
      \coordinate (A) at (-3.9,0); \coordinate (B) at (-2.0,1.8);
      \coordinate (C) at (1.0,1.8); \coordinate (D) at (3.9,0);
      \coordinate (E) at (1.5,-1.8); \coordinate (F) at (-2.2,-1.8);
      \draw[geom side] (A)--(B)--(C)--(D)--(E)--(F)--cycle;
      \draw[geom diagonal] (A)--(C) (A)--(D) (A)--(E);
      \foreach \P/\pos in {A/left,B/above left,C/above,D/right,E/below,F/below}{
        \node[geom point] at (\P) {}; \node[geom label,\pos=3pt] at (\P) {$\P$};}
      \node[geom note] at (0,-2.45) {Из $A$ выходят три диагонали: $AC$, $AD$, $AE$};
    \end{tikzpicture}
  \end{center}
}
